% AI-written, AI cross-reviewed; no human sign-off; not a claim that RH is proved
% RH Notes blueprint source excerpts for Z1.2
% Verbatim excerpts in the notation of the parent blueprint. The labelled editions and source ranges appear below.
% See the unit YAML for the current proof status. Earlier-round passages retain their original wording.

% Start of atlas edition annotation. The following block is outside the original excerpt.
\paragraph{Status of named inputs in this reading.}
\noindent\nolinkurl{UpperTheory.V.Loc.LocalPublishedTheorems.bk} --- record field (\path{Development/AutomorphicArithmetic/AutomorphicVarieties/PadicHodgeTheory/LocalPublishedTheorems.lean:19}); not compared.\par
% End of atlas edition annotation.

\section*{Z1.2. Ordinary adjoint rigidity}

% Parent source SHA-256: 1af64922b98ed2310744411b1f18b25cc426c97220646321471836e1a3a409d6
% blueprint/ch_P.tex lines 12645 through 12652; r6:review:Z1.2
\paragraph{Z1.2. Local and global proof boundaries.}\label{r6:review:Z1.2} \emph{Original argument on p.~\pageref{sec:P-r5-z12}.} 
The declared root takes ordinary vanishing and the local published record as premises. Its \nolinkurl{OrdZero} input uses a form over $c.O$ and its \nolinkurl{AdZero} target uses the form condition over $\mathbf Z_\ell$. The two multiplier characters agree after extension to $\overline{\mathbf Q}_\ell$ by the scalar tower identities. Source: \nolinkurl{adjointLiftingTheorem_of_locPrints_of_core}, \path{Development/AutomorphicArithmetic/AutomorphicVarieties/ModularityArgument/ThresholdArguments/DenseThresholdCases/AdLOf.lean:104--150}. The following inputs remain \emph{proved in Lean, not reproduced}.
\begin{itemize}
\item For the local step, transport the decomposition group, its inertia and cyclotomic character through an open embedding, then transport the finite period primitive (\nolinkurl{Loc.decEmb}, \path{Development/AutomorphicArithmetic/AutomorphicVarieties/PadicHodgeTheory/DecEmb.lean:71}; \nolinkurl{Loc.bkAbove}, \path{Development/AutomorphicArithmetic/AutomorphicVarieties/PadicHodgeTheory/BkAbove.lean:47}). Crystallinity of the ordinary representation on the small subgroup uses the decreasing weights and exclusions of cyclotomic links (\nolinkurl{Loc.prC}, \path{Development/AutomorphicArithmetic/AutomorphicVarieties/PadicHodgeTheory/SnocZeroLastCoordinateLast.lean:345--352}); passage to its endomorphisms uses internal Hom (\nolinkurl{Loc.isCrys_isDR_internalHom}, \path{Development/AutomorphicArithmetic/AutomorphicVarieties/PadicHodgeTheory/AdmHom.lean:108}). Restriction of finite cocycles to the stable corner is \nolinkurl{Th.z1fSub}, \path{Development/AutomorphicArithmetic/AutomorphicVarieties/TheoremsMuLinked/ContinuousSubActions/ContActionSub.lean:83}; the negative- and zero-weight scalar steps are \nolinkurl{Loc.z1fChar}, \path{Development/AutomorphicArithmetic/AutomorphicVarieties/PadicHodgeTheory/ExistsNeOne.lean:63}\ and \nolinkurl{Loc.z1fUnr}, \path{Development/AutomorphicArithmetic/AutomorphicVarieties/PadicHodgeTheory/ExistsIntBasis.lean:370}. The coefficient descent is \nolinkurl{OrdinaryBlock.desc}, \path{Development/AutomorphicArithmetic/AutomorphicVarieties/TheoremsMuLinked/GaloisRepresentationTypes/CornerStable.lean:273--305}.
\item For ordinary vanishing the precision-dependent auxiliary places require \nolinkurl{Pa.arPl}, \path{Development/AutomorphicArithmetic/AutomorphicVarieties/PatchingModules/AutomorphicDetection/ArPl.lean:13}. The uniform constant in the global upper bound uses the finite-level and dual Selmer estimates (\nolinkurl{Th.oneLevelO}, \path{Development/AutomorphicArithmetic/AutomorphicVarieties/TheoremsMuLinked/CyclicReductionCohomology/OneLevelO.lean:52}; \nolinkurl{Th.b1absO}, \path{Development/AutomorphicArithmetic/AutomorphicVarieties/TheoremsMuLinked/CyclicReductionCohomology/ResidualRepOne.lean:44}). The lower bound needs the Hida ideal, regularity and support of an operator with nonzero specialization (\nolinkurl{Hi.exists_hecke_ideal}, \path{Development/AutomorphicArithmetic/AutomorphicVarieties/ModularityArgument/ThresholdArguments/DenseThresholdCases/ExistsHeckeIdeal.lean:70}; \nolinkurl{Th.Nm.regular_arithmetic_hida_of_literal_automorphy}, \path{Development/AutomorphicArithmetic/AutomorphicVarieties/ModularityArgument/ThresholdArguments/DenseThresholdCases/NmRegArith.lean:109}; \nolinkurl{Th.has_operator_statement}, \path{Development/AutomorphicArithmetic/AutomorphicVarieties/ModularityArgument/ThresholdArguments/DenseThresholdCases/THasOp.lean:18}). The weight algebra, enlarged determinant ring, trace tests and integral model are further inputs to this count. Their assembly occurs in \nolinkurl{Th.hida_x_of_printed_hida}, \path{Development/AutomorphicArithmetic/AutomorphicVarieties/ModularityArgument/ThresholdArguments/DenseThresholdCases/THidaX.lean:17}\ and \nolinkurl{Th.core_alt_bis}, \path{Development/AutomorphicArithmetic/AutomorphicVarieties/ModularityArgument/ThresholdArguments/DenseThresholdCases/ThFinal.lean:17--18}.
\item For unlinking, the local unramified direction and cochain (\nolinkurl{Z.J.cupUnr}, \path{Development/AutomorphicArithmetic/AutomorphicVarieties/ModularityArgument/LinkedUnramifiedChains/CrystallineCardinalities/CupUnr.lean:101}), ordinary and exceptional local counts (\nolinkurl{Z.J.locLow}, \path{Development/AutomorphicArithmetic/AutomorphicVarieties/ModularityArgument/LinkedUnramifiedChains/CrystallineCardinalities/LocLow.lean:134}; \nolinkurl{Z.J.locHigh}, \path{Development/AutomorphicArithmetic/AutomorphicVarieties/ModularityArgument/LinkedUnramifiedChains/CrystallineCardinalities/LocHigh.lean:18}), and the associated-prime/derivation obstruction (\nolinkurl{V.J.derZero}, \path{Development/AutomorphicArithmetic/AutomorphicVarieties/ModularityArgument/LinkedUnramifiedChains/DerivativeDivisibility/DerZero.lean:53}) are substantive. A proof still needs the normalization and global duality bounds, purity and simple-root selection on the same component, and the two signed monomials. The last flag argument uses splitting to lift local elements and preservation of potential crystallinity under restriction, isomorphism and duality (\nolinkurl{move_all, unlEnd3E}, \path{Development/AutomorphicArithmetic/AutomorphicVarieties/ModularityArgument/LinkedUnramifiedChains/LogarithmicNormBounds/BlockValuationRep.lean:150--215,234--346}).
\end{itemize}
\emph{Hypothesis, discharged at the application.} The actual call to \nolinkurl{Z.J.linkedUnlE_alt}, \path{Development/AutomorphicArithmetic/AutomorphicVarieties/ModularityArgument/ThresholdArguments/DenseThresholdCases/LinkedUnlE.lean:43--48} supplies Artin reciprocity, local class field theory, unitary modularity, \nolinkurl{modularity_n_of_printed_mult_comp}, the level and weight algebras, the universal determinant ring, both seed ranks, twist and algebraic character comparisons, local/global cohomological comparisons, the tangent input and ordinary vanishing. Thus the published local record, the derived finite-to-ordinary comparison and the separate integral vanishing and unlinking arguments have identifiable, different proof obligations.

% Parent source SHA-256: 509893882390ca545e5542b088ec432ed3c1ed009152eeca67e7fc44894dd956
% blueprint/ch_P_r11.tex lines 32 through 432; r11:adj
\subsection{Z1.2. Three different rigidity assertions}
\label{r11:adj}
Throughout this subsection file suffixes are relative to
\path{Development/AutomorphicArithmetic/AutomorphicVarieties/}.
Write $D$ for its directory
\path{ModularityArgument/ThresholdArguments/DenseThresholdCases/},
$C$ for \path{TheoremsMuLinked/CyclicReductionCohomology/}, and $J$ for
\path{ModularityArgument/LinkedUnramifiedChains/}.

\paragraph{The declared root.}
Let $L$ be the CM field in the Hida setting, $V.S'=s.S$, with no primitive
$\ell$th root in $L$, $c$ large, $2\le n\le4$, $r.\mathrm{V0ok}$,
$B$ the universal ordinary determinant ring, and $x$ the arithmetic point
with locally generic automorphic realization $\rho$.
In particular the given flag has characters $\psi_{v,i}$, and the freeness,
irreducibility and separability hypotheses are those in
\nolinkurl{AutomorphicPointLocallyGeneric}, $D$\path{AutomorphicPointGalois.lean:18--39}.
For each $v\mid\ell$, assume there are strictly decreasing integers $m_i$
and a single $k_v>0$ such that
$(\psi_{v,i}/\psi_{v,j}\,\epsilon^{m_j-m_i})^{k_v}=1$ on inertia
for all $i,j$. For every allowed polarization and polarized adjoint action
$R$, let $z$ be a continuous cocycle whose pullback to $G_{L^+}$ is
Bloch--Kato finite. The conclusion is one $X$ with $z(g)=R(g)X-X$ for
\emph{every} $g$. The declared implication assumes
\nolinkurl{LocalPublishedTheorems} and \nolinkurl{DenseThresholdMainCase}.
The latter supplies this conclusion when $z$ vanishes on $r.I_p(u)$ and,
locally at each $v$, subtracting one coboundary in an ordinary basis
makes $z$ upper triangular everywhere and its diagonal zero on inertia.
The polarization character obtained over $c.O$ agrees with that over
$\mathbf Z_\ell$ after mapping into $\overline{\mathbf Q}_\ell$ by transitivity
of the coefficient maps.
\nolinkurl{OrdAt, OrdCoc, Par, OrdZero, AdZero},
\path{TheoremsMuLinked/OrdZero.lean:22--34,56--72};
\nolinkurl{adjointLiftingTheorem_of_locPrints_of_core},
$D$\path{AdLOf.lean:104--150}.

Away from $\ell$ take the primitive at the place provided by the finite
condition. Upon conjugating the place by $\gamma$, the primitive becomes
$R(\gamma)X-z(\gamma)$: expand $z(g\gamma)=z(\gamma(\gamma^{-1}g\gamma))$.
On the specified inertia the determinant relations give
$\mathrm{tr}(\rho(h)(\rho(g)-1))=0$ for every $h$. The matrices
$\rho(h)$ span the full matrix algebra by irreducibility over the
algebraically closed coefficient field. Testing against matrix units gives
$\rho(g)=1$, hence $z(g)=0$ there. Above $\ell$ two embeddings over the
base differ by an automorphism of its algebraic closure. The same
calculation transports a finite primitive $x$ to
$(R(\gamma)\otimes1)x-z(\gamma)\otimes1$: the conjugation intertwines
the two local actions on simple tensors and hence on their sums.
The open embedding matching inertia and cyclotomic characters and the
finite-to-ordinary implication then give the second ordinary condition.
This is the deduction at the root conditional on those local comparisons.
\nolinkurl{adL_conj, adL_emb, adL_above, adL_away, adL_desc, adL_triv},
$D$\path{AdLOf.lean:10--103};
\nolinkurl{bkOrd}, \path{TheoremsMuLinked/GaloisRepresentationTypes/CornerStable.lean:373--436}.

\paragraph{Finite local classes and the ordinary condition.}

For an open embedding $\phi:G\hookrightarrow G_{\mathbf Q_\ell}$ write
$dX(g)=R(g)X-X$. A continuous cocycle $y$ is finite if there is one
$x\in V\otimes_{\mathbf Q_\ell}B_{\rm cris}$ such that
$y(g)\otimes1=gx-x$ for every $g$. The ordinary conclusion is a single
$X$ making $y-dX$ upper triangular at every $g$ in a flag basis, with
zero diagonal on inertia (\nolinkurl{BkOrd}, \path{FreeFlag.lean:33--40};
\nolinkurl{OrdinaryBlock.fin}, \path{CornerStable.lean:122--129}).

\paragraph{Away from $\ell$.}
For $g\in r.I_p(u)$ the determinant relation gives
$\operatorname{tr}(\rho(h)(\rho(g)-1))=0$ for \emph{every} $h$.
Indeed the relation tests $[g]-[1]$ against $[h]$ before specialization,
and the trace of the specialized determinant is the matrix trace.
Irreducibility gives full matrix span at the given point. Linearity
therefore permits $\rho(h)$ to be replaced by $E_{ji}$, which reads
off the $(i,j)$ entry of $\rho(g)-1$. Thus $\rho(g)=1$, and its
adjoint action fixes every primitive on this inertia. Sources:
\nolinkurl{Ws.Eb.relKill}, \path{EigenbasisLifts/RelKill.lean:13--46};
\nolinkurl{Ws.Eb.trPlus}, \path{EigenbasisLifts/TrPlus.lean:10--42};
\nolinkurl{adL_triv}, \path{AdLOf.lean:98--122}.

\paragraph{Transport of a period primitive.}
For $j:\Omega\hookrightarrow\overline{\mathbf Q}_\ell$ and number field
$F$, let $K=\mathbf Q_\ell(j(F))$; adjoining images of a rational
basis suffices. Restriction sends $G_K$ to the decomposition group of
$A=\{x:\lVert jx\rVert\le1\}$ by
$j(\operatorname{res}\sigma(x))=\sigma(jx)$. Normality gives the
restriction, preservation of the norm puts it in the stabilizer, and
injectivity of $j$ gives the product law. For continuity at $1$ take
a finite extension fixed by a basic Krull neighbourhood. Stabilizers
of the images of a finite basis have open intersection mapping into the
neighbourhood. Restriction is injective by density of $j(\Omega)$.
Conversely $g\in\operatorname{Dec}(A)$ preserves the balls
$\lVert jz\rVert\le|\ell|^m$ by the valuation-ring test on $z/\ell^m$.
It extends to the completion, fixes $\mathbf Q_\ell$ by density of
$\mathbf Q$, and restricts back to the normal algebraic closure. It fixes
$j(F)$. Compactness and Hausdorffness give a homeomorphism. To compare
inertia, approximate each local integral element by $jx$ at distance
$<1$ and apply the ultrametric inequality. To compare cyclotomic
characters, compare actions on a primitive $\ell^m$th root for every $m$.
Sources: \nolinkurl{CovariantDegrees.j_resE},
\nolinkurl{CovariantDegrees.rHom}, \nolinkurl{CovariantDegrees.continuous_rHom},
\path{IntegralCasts/CapelliIdentities/CovariantDegrees.lean:1232--1249},
\path{CapelliIdentities/CovariantDegrees.lean:4809--6170};
\nolinkurl{LB.denseRange_emb}, \path{NDLocA.lean:12--45};
\nolinkurl{LB.KwF}, \path{NDLocB.lean:14--24};
\nolinkurl{LB.rHom_injective}, \nolinkurl{LB.exists_ext},
\nolinkurl{LB.rHom_surjective}, \path{NDLocC.lean:87--184};
\nolinkurl{LB.eDec}, \path{NDLocD.lean:21--35};
\nolinkurl{de_ine}, \nolinkurl{de_cyc}, \nolinkurl{Loc.decEmb},
\path{PadicHodgeTheory/DecEmb.lean:12--74}.

The group of compatible global/local pairs $\operatorname{decJ}(j)$
projects onto $G_K$ with inverse pairing a local automorphism and its
restriction. Existence of a period primitive persists on restriction to
an open subgroup. Conversely, if $f(h)=hx-x$ on a finite-index subgroup
$H$, then $u=f-dx$ satisfies $u(gh)=u(g)$ for $h\in H$. On left
cosets $u(gq)=u(g)+gu(q)$. For $S=\sum_{q\in G/H}u(q)$ this gives
$[G:H]u(g)=S-gS$, so $x-S/[G:H]$ is a global primitive. Changing
the embedding by $\gamma$ replaces the primitive by
$(R(\gamma)\otimes1)x-c(\gamma)\otimes1$, by the cocycle identity.
Sources: \nolinkurl{Af.decJ}, \path{BlochKatoSelmerVanishes.lean:119--128};
\nolinkurl{Loc.decJFacts}, \path{DecJFacts.lean:5--50};
\nolinkurl{bkAbove_avg}, \nolinkurl{Loc.bkAbove}, \path{BkAbove.lean:5--79};
\nolinkurl{adL_above}, \path{AdLOf.lean:27--57}.

\paragraph{Finite coefficient descent.}
For compact Hausdorff $G$ and a continuous matrix representation over
$\overline{\mathbf Q}_\ell$, take a countable family of finite extensions
cofinal on finite sets of coefficients. To construct it, approximate an
algebraic coefficient by an element of the embedded $\overline{\mathbf Q}$
more closely than its distances to its other conjugates. Krasner's lemma
puts it in the field of the approximation; use finite sets of approximants.
The closed subgroups on which entries of both $r(g)$ and $r(g^{-1})$ lie
in the respective field cover $G$. Baire gives one with interior; it is
open and has finite index. Adjoining entries of coset representatives
gives a finite extension $E_0$ containing all representation entries.
Sources: \nolinkurl{CompactRepField.exists_mem_adjoin_simple},
\nolinkurl{CompactRepField.exists_countable_family},
\nolinkurl{CompactRepField.exists_finiteDimensional_entries_mem},
\path{CompactGroupRepDefinedOverFiniteExtension.lean:14--119}.

Write the period primitive as a finite sum $x=\sum_t A_t\otimes b_t$.
Adjoin all entries of the $A_t$ to $E_0$ to obtain $E$, and choose an
$E$-linear retraction $p:\overline{\mathbf Q}_\ell\to E$ of the
inclusion $i$. Entrywise $p$ commutes with conjugation by $r(g)$.
Thus $(p\otimes1)x$ is a primitive for $p y$, and $(i\otimes1)$
takes it back to $x$. It follows that
$i p y(g)\otimes1=y(g)\otimes1$. Tensoring with $1$ is injective:
the nonzero $\mathbf Q_\ell$-algebra $B_{\rm cris}$ admits a linear
functional taking $1$ to $1$. So $y$ itself descends, and continuity
of $p y$ follows from the induced topology and $i p y=y$. All
diagonal characters descend as entries. Sources:
\nolinkurl{OrdinaryBlock.tmul_one_inj}, \nolinkurl{OrdinaryBlock.desc},
\nolinkurl{OrdinaryBlock.bar}, \path{CornerStable.lean:196--201},
\path{CornerStable.lean:209--345}.

\paragraph{Rank-one and extension comparisons.}
Write $H_f^1=Z_f^1/B^1$. Fix finite $E/\mathbf Q_\ell$ and a continuous
character $\kappa$ of weight $w<0$ on an open inertia subgroup.
The cyclotomic character has an element of infinite order on this
inertia, so $\kappa$ is nontrivial. Multiplication by
$\kappa(g)-1\ne0$ at one $g$ shows $H^0=0$ and $\dim_E B^1=1$.
The rank-one filtration and the local dimension formula give
\[
 \operatorname{Fil}^0D_{\rm dR}=D_{\rm dR},\qquad
 \dim_E H_f^1+\dim_E\operatorname{Fil}^0D_{\rm dR}
 =\dim_E H^0+[G_{\mathbf Q_\ell}:\phi(G)] .
\]
De Rham admissibility and restriction of scalars give
$\dim_E D_{\rm dR}=[G_{\mathbf Q_\ell}:\phi(G)]$. Hence $H_f^1=0$:
every finite scalar cocycle is $g\mapsto(\kappa(g)-1)m$.
Sources: \nolinkurl{Loc.z1fChar}, \path{ExistsNeOne.lean:31--106};
\nolinkurl{Loc.charFil}, \path{CharFil.lean:193--232};
\nolinkurl{Loc.bkDim}, \path{CocycleTransportH1fSurjective.lean:133--173}.

For a continuous cocycle $c$ form $E_c=V\oplus E$ with
$g(v,a)=(gv+a c(g),a)$. A $G$-equivariant injection $f:V\to W$
induces an injection $E_c\to E_{fc}$ by $(v,a)\mapsto(fv,a)$.
If $W$ is crystalline on an open normal subgroup $U$ and $fc$ is
finite, the extension criterion gives crystallinity of $E_{fc}$ on $U$.
Closure under subrepresentations gives this for $V$ and $E_c$; the
criterion gives finiteness of $c$. Applied to a stable matrix corner
this permits entry evaluation in its rank-one quotient. On the trivial
diagonal character, a finite cocycle vanishes on inertia: its extension
is crystalline, the filtration has steps $\mathrm{Fil}^0=D$,
$\mathrm{Fil}^1=0$, and trivial inertia on the extension evaluated
at $(0,1)$ forces $c(g)=0$. Sources: \nolinkurl{Th.z1fSub},
\path{ContActionSub.lean:58--100}; \nolinkurl{OrdinaryBlock.entry},
\path{CornerStable.lean:72--93}; \nolinkurl{Loc.z1fUnr},
\path{PadicHodgeTheory/ExistsIntBasis.lean:303--375};
\nolinkurl{Loc.crWt0}, \path{WtMatFix.lean:204--265}.

The triangular-to-crystalline step here takes $A=E$ and $\pi=1_E$ on
$U$. Write the action in blocks $(M_0,b;0,k)$ with $k$ the last
character. By induction the top block is crystalline. The extension
cocycle is $c(g)=b(g)/k(g)$ for $V_0\otimes k^{-1}$, by block
multiplication. The twist is crystalline, and its flag factors have
weights $w_i-w_{\rm last}>0$. The pairwise hypothesis excludes the
cyclotomic character among them. A nonzero equivariant functional to
the cyclotomic line would descend nontrivially from the first flag
quotient on which it is nonzero. Thus the twisted dual invariants
vanish. The positive-weight comparison makes $c$ finite, and the
extension criterion completes the induction. For a finite algebra over
$E$, powers of the nilpotent kernel of $\pi:A\to E$ refine each flag
quotient into residual scalar factors. For the full crystalline column
space $V$ restricted to $\mathbf Q_\ell$, evaluation embeds the matrix
space into $V^*\otimes_{\mathbf Q_\ell}V$. Stability under dual, tensor,
and subrepresentation gives its crystallinity. Sources:
\nolinkurl{Loc.prC}, \nolinkurl{DualInvariantFiltration.key},
\path{SnocZeroLastCoordinateLast.lean:53--146},
\path{SnocZeroLastCoordinateLast.lean:217--352};
\nolinkurl{OrdinaryBlock.crOn_ad}, \path{CornerStable.lean:96--120};
\nolinkurl{Loc.isCrys_isDR_internalHom}, \path{AdmHom.lean:18--27},
\path{AdmHom.lean:77--86}, \path{AdmHom.lean:108--120}.

\emph{Proof boundary.} The comparisons used in these deductions are
proved in Lean; their foundational period-ring and cohomological
arguments are not reproduced here. With prefix \nolinkurl{UpperTheory.V.Loc},
the seven comparison roots and their inclusive core line counts are
\nolinkurl{crExt} (crystallinity of a rank-one quotient extension
equivalent to finiteness of its twisted cocycle, 18,598,
\path{MatrixExtendByZeroApply.lean:201--273}),
\nolinkurl{prF} ($Z_f^1=Z^1$ for a continuous positive representation
with zero cyclotomically twisted dual invariants, 20,759,
\path{PrF.lean:7--23}), \nolinkurl{bkDim} (11,421),
\nolinkurl{charFil} (16,908), \nolinkurl{crWt0} (16,171),
\nolinkurl{isCrys_isDR_sub_quot} (admissibility under equivariant
injections and surjections, 13,208, \path{AdmSub.lean:105--129}),
and \nolinkurl{isCrys_isDR_internalHom} (16,296).
The union of these seven inclusive cones has 26,405 core lines.

\paragraph{Ordinary tangent count.}
The precise upper bound concerns a trace image. For a fixed Taylor--Wiles
system $T$ with $q$ auxiliary places, integral point $y$, total local degree
$d$, $Q=|O/\mathfrak m|$ and $b=nd+nq$, write
\[
 S_{m,N}=\{g\mapsto\mathrm{tr}(z(g)\rho_{m,N}(g)):\ 
                 z\in\mathrm{Sel}_{m,N}\}.
\]
The set on the right is taken over locally constant polarized cocycles
vanishing on marked inertia and satisfying the local ordinary conditions.
The conclusion of \nolinkurl{b1absO} is
$\exists C_*,\ \forall m,\ \forall N\gg_m0,\quad
 S_{m,N}\text{ finite},\ |S_{m,N}|\le Q^{mb+C_*}$.
Its hypotheses include $\ell\ne2$, the system with degree $d$ and
subgroup $D_0$, \nolinkurl{PhiCond}, a nontrivial local cyclotomic value
at each $v$, triviality of $y.\rho$ on the prescribed \nolinkurl{PL}
intersection, $y.\mathrm{N4}(D_0)$, distinct split auxiliary eigenvalues
and the integral trace detection condition \nolinkurl{DetO}.
\nolinkurl{Pt.Sel, Pt.trFun, Pt.B1},
\path{TheoremsMuLinked/OrdinaryHeckeAlgebraIntegral.lean:52--66};
\nolinkurl{b1absO}, $C$\path{ResidualRepOne.lean:44--82};
\nolinkurl{DetO}, $C$\path{DualPair.lean:113}.

Here is the step that makes the precision quantifier work.
Put $t={n\choose2}$ and let $d_v$ sum to $d$. Take the global invariant
and dual bounds $c_1,c_2$, the distinguished auxiliary bound $c_0$, the
ordinary local constants $C_v$ and the split auxiliary constants $A_j$
\emph{before} choosing $m$. After $m$ is fixed intersect the two eventual
level conditions for dual cocycles and the auxiliary Frobenius tests.
At such a level the one-level product inequality bounds
$|S_{m,N}|Q^{u}$, where $u=m(n+t)\#\{v\}+mtd$, by $Q$ to the power
\[
 c_1+c_2+\sum_v\{m(n+t)(1+d_v)+C_v\}+c_0+\sum_{j=1}^q(mn+A_j).
\]
Subtracting $u$ gives $mb+C_*$ with
$C_*=c_1+c_2+c_0+\sum C_v+\sum A_j$. Cancellation is valid since $Q>0$.
The one-level argument uses the ordinary local subgroup generated by
coboundaries and the triangular subgroup, the cocycles vanishing on
the marked inertia at tame places, and the full cocycles at auxiliary
places. These lie between coboundaries and cocycles, so the given global
duality formula applies. Inclusion of the trace image is what permits the
last cardinality bound. \nolinkurl{oneLevelO}, $C$\path{OneLevelO.lean:52--134};
\nolinkurl{sum_exp, b1absO}, $C$\path{ResidualRepOne.lean:25--33,44--129}.

\emph{Statement only (project proof in Lean).}
In particular the dual estimate still needed in this reduction asserts,
under the odd sign at conjugation, \nolinkurl{PhiCond}, the prescribed
\nolinkurl{PL} triviality and \nolinkurl{DetO}, that there is a fixed
$c_2$ such that, for each $m$ and all sufficiently large $N$, any finite
space of locally constant cocycles of the specified twisted adjoint
action $R'$ has
\[
 |\{z\in Z^1_{\rm lc}(R'):\ 
     \forall j\ \exists Y_j,\ z(s)=R'(s)Y_j-Y_j
                  \ (s=\phi_{N,j},\tau_{N,j})\}|
 \le |B^1(R')|Q^{c_2}.
\]
The same $Y_j$ is required at the two generators. Its source first descends
a cocycle normalized at conjugation and then bounds integral twisted
cocycles with vanishing traces against all powers of the test matrices.
The latter uniform bound and the construction of the weight and auxiliary
derivations on the same Hecke level have no complete hand proof here.
\nolinkurl{dualNO}, $C$\path{DualNO.lean:50--55,95--144}
(cone 2,410). The containing \nolinkurl{b1absO} cone is 4,280.
The separate application \nolinkurl{core_alt_bis}, $D$\path{ThFinal.lean:17},
still requires the ordinary vanishing assertion for all the points specified
above (cone 375,722). Its two-sided count is applied in
\nolinkurl{core_o_of_literal_automorphy_n}, $D$\path{TCoreO.lean:129--159},
and \nolinkurl{ordinary_vanishing_two_of_printed_hida},
$D$\path{TOrdVan.lean:159--192}.

\paragraph{Unlinked flag: two deductions at the exceptional place.}
First suppose an adjacent cyclotomic extension supplies
\nolinkurl{UnitStep} in a triangular rank-three representation over a
coefficient DVR $O$. Write $A=O/(\ell^k)$, $a=|A|$, $k>0$, and $d$ for
the local degree. In the space $Z_F$ of upper triangular ordinary
cocycles let $W$ be the subgroup whose two adjacent diagonal functions
are equal everywhere. Suppose $|W|\ge a^{6+3d}$ and at every $k$ a
cocycle $y\in Z_F$ has adjacent diagonal difference $\bar c$ at some
$s_0$, for a fixed $0\ne c\in O$. For each
$b\in(\bar c)\subset A$ choose $r_b\bar c=b$. Then
$(w,b)\mapsto w+r_by$ is injective: the adjacent diagonal difference
at $s_0$ recovers $b$, after which subtraction recovers $w$.
Since $O/(c)$ surjects onto $A/(\bar c)$, we obtain the uniform estimate
\[
 a^{7+3d}\le |O/(c)|\,|Z_F|.
\]
This is the extra local power \emph{before} adding the unrestricted
matrix coboundaries. The two supplier statements here are the lower bound
on $W$ and existence of $c,y$ from the local Euler and generation data,
triangularity, $(\ell:O)\ne0$ and \nolinkurl{UnitStep}. The group is
compact and the coefficients carry the specified $\ell$-adic structure.
\nolinkurl{locHigh}, $J$\path{CrystallineCardinalities/LocHigh.lean:18--87}
(cone 1,109); \nolinkurl{sel_cut,selLow},
$J$\path{CrystallineCardinalities/SelLow.lean:44--52,60--98}.

The cochain used for \nolinkurl{UnitStep} has a prescribed sign. For an
open embedding of the local group in $G_{\mathbf Q_\ell}$, finite
$E/\mathbf Q_\ell$ and a finite cyclotomic cocycle $\kappa$, restrict the
unramified additive degree $d_0$ with $d_0(\mathrm{Frob})=1$. Some
positive power of Frobenius belongs to the open image, so $d_0$ there is
nonzero. For a $\mathbf Q_\ell$-basis $e_i$ of $E$, the trivial-action
cocycles $d_0e_i$ are finite by the unramified comparison. Pair each with
$\kappa$ using the dual-coordinate pairing $E(1)\times E\to\mathbf Q_\ell(1)$.
The finite-class orthogonality in the assumed Bloch--Kato record supplies
continuous $f_i$ with
$\epsilon(s)f_i(t)-f_i(st)+f_i(s)=d_0(t)\kappa_i(s)$.
Hence $u=d_0$ in $E$ and $\beta=\sum_i f_ie_i$ satisfy
$\beta(st)=\beta(s)+\epsilon(s)\beta(t)-\kappa(s)u(t)$.
This deduction uses the comparison of the finite local conditions with
those in the record and the unramified criterion, including on the
open subgroup. \nolinkurl{exists_deg,gk_case,cupUnr},
$J$\path{CrystallineCardinalities/CupUnr.lean:13--22,43--97,101--135}
(cone 24,143).

\paragraph{From counting to two integral tangent vectors.}
Here the passage from finite levels to characteristic zero admits a
self-contained module argument. Assume $O$ a coefficient DVR with
$(\ell:O)\ne0$, $R$ a Noetherian commutative $O$-algebra, $\xi:R\to O$
an $O$-algebra point and $T\subset R$. Suppose the number of $O$-linear
$\xi$-derivations to $A_k=O/(\ell^k)$ killing $T$ is at least
$|A_k|^2/C$ for all $k>0$, with $C>0$ fixed. Then there are two
$O$-linearly independent such derivations to $O$.
Form $M=R/\langle T,ab-\xi(a)b-\xi(b)a\rangle_O$ as an $O$-module.
Derivations with these conditions are exactly the linear maps from $M$.
In $M$ the class of $1$ vanishes (take $a=b=1$), and so does the square
of $\ker\xi$ as an ideal. The finite ideal generators of $\ker\xi$
therefore span $M$ over $O$: replace each coefficient $r$ by $\xi(r)$,
since their difference lies in the kernel. Decompose $M$ over the DVR
as a free summand of rank $s$ and cyclic torsion summands $O/(a_i)$.
In $A_k$, multiplication by $a_i$ has kernel and cokernel of equal order.
The cokernel is a quotient of $O/(a_i)$, so the possible images of the
torsion generators number at most $C'=\prod_i|O/(a_i)|$ independently of
$k$. If $s\le1$, there are at most $C'|A_k|$ maps. Since $|A_k|$ tends
to infinity this contradicts the lower bound. Projection on two different
free coordinates gives the independent integral derivations. The proof
takes no inverse limit of choices. \nolinkurl{tors_card,hom_alt,rel_fin,rankTwo},
$J$\path{CrystallineCardinalities/RankTwo.lean:9--96,98--165,177--229}.

\paragraph{The flag obligation that remains.}
On an even level, for each linked frame with \nolinkurl{Ok}, \nolinkurl{Big}
and its \nolinkurl{SeedClause}, the assertion \nolinkurl{UnlPt} quantifies
over the specified integral rank-three realizations with prescribed inertia,
their compatible Galois realization irreducible on all open subgroups, and
each frame place $v$ at which it is potentially crystalline. It requires
\emph{for every ordered pair} $i,j$ an element $g\in\Phi.G_v$ such that
\[
 \psi_{v,i}(g)\ne\psi_{v,j}(g)\,\epsilon(k.j_v(g))
 \quad\text{after the specified coefficient embedding.}
\]
The independent derivations constructed from a hypothetical link must
annihilate the entire weight-zero algebra on the \emph{same} component
before the type-I grading gives a contradiction. In particular the
annihilator, simple-root selection and the local-to-global lower count at
that point remain project inputs. For the final comparison of flags the
matching argument itself is elementary. In one triangular basis with
diagonal $D_b$, suppose none of the frame characters $\psi_j$ agrees
identically with $D_b$. Choose $g_j$ giving a mismatch for each $j$.
The $bb$-entry of the ordered product
$\prod_j(T(g_j)-\psi_j(g_j)I)$ is the product of three nonzero scalars.
The defining ordinary product is zero, a contradiction. This gives a
single matching character per row, sufficient for the application, even
without a fixed permutation in advance. Transport to a frame place and
restriction of potential crystallinity are additional hypotheses of the
comparison in the endpoint proof.
\nolinkurl{UnlPt}, $J$\path{WeightInfinitesimalCharacters/AdmissibleUnitaryThreePoint.lean:20--27};
\nolinkurl{LinkedUnlE,UnlEndLogNorm}, $J$\path{LogarithmicNormBounds/LinkedUnlEOf.lean:12--30};
\nolinkurl{flag_match,unlEnd3E}, $J$\path{LogarithmicNormBounds/BlockValuationRep.lean:27--63,234--346};
\nolinkurl{linkedUnlE_alt}, $D$\path{LinkedUnlE.lean:43--48}
(cone 1,111,243). This cone and the ordinary tangent input are at callers
outside the declared root's 38,519-line cone. The counts are inclusive
core declaration spans and overlap.

